%%%%%%%%%%%%%%%%%%%%%%%%%%%%%%%%%%%
% Main Text
%%%%%%%%%%%%%%%%%%%%%%%%%%%%%%%%%%%

\section{Introduction}\label{sec:intro}

Asian Americans represent the fastest-growing racial group in the United States, yet their experiences with discrimination and identity formation remain underexplored in economic research.\footnote{The 2020 Census counted more than 20 million Asian Americans---6.4 percent of the population—nearly double the number counted two decades earlier. The Asian American population numbers are based on the author's calculations from the Current Population Survey and US Census data \autocite{floodsarahIntegratedPublicUse2021a}.} Unlike other minority groups, Asian Americans occupy a distinctive position in America's racial hierarchy---simultaneously experiencing discrimination and being labeled as ``perpetual foreigners'' while also being characterized through the ``model minority'' stereotype \autocite{foukaImmigrantsAmericansRace2022}. This dual status creates complex incentives around racial identity choices that fundamentally differ from other groups' experiences. Understanding how both anti-Asian bias and socioeconomic status influence racial identity decisions is crucial, as the selection could systematically distort the true enumeration of the Asian population in the US, the empirical assessments of Asian American outcomes and the true extent of discrimination they face.

An extensive literature has documented how Asians and Whites compare across multiple outcomes, yet the role of identity selection in shaping the Asian-White gaps remains understudied \autocite{chiswick1983analysis, duleep2012economic, hilger2016upward, arabsheibani2010asian}. The challenge lies in defining and measuring racial identity, particularly when individuals possess agency in how they racially self-identify. If reporting Asian racial identity represents a strategic choice influenced by local discrimination, measured gaps may systematically vary across geographic contexts in ways that previous research has not fully explored.

Various contextual factors, including anti-Asian sentiment and stereotypes, can influence how individuals navigate their racial identity choices. The COVID-19 pandemic has brought renewed attention to how external hostility shapes Asian American experiences \autocite{gover2020anti}. In this paper, I examine the determinants of Asian racial identity reporting and analyze how Asian Americans strategically select between Asian and White racial identities. Specifically, I investigate how anti-Asian bias, education, and family income shape decisions to identify racially as Asian American. I also break down the analysis by generation and family structure (interracial versus endogamous parents) and examine what racial identities individuals with objective Asian ancestry report (Asian only, White only, multiracial, etc.).

I explore how individual characteristics and societal attitudes toward Asian Americans influence racial identity reporting. I use identity and ancestry data from the Current Population Survey (CPS) combined with measures of anti-Asian bias derived from Harvard's Project Implicit Association Test (IAT), the American National Election Studies (ANES), and hate crimes targeting Asian Americans.\footnote{The IAT data comes from Harvard's Project Implicit \autocite{greenwaldMeasuringIndividualDifferences1998}. Implicit bias measures have gained prominence in economics, with IAT scores correlating with economic outcomes \autocite{chettyRaceEconomicOpportunity2020,gloverDiscriminationSelfFulfillingProphecy2017}, voting patterns \autocite{friesePredictingVotingBehavior2007}, and health disparities \autocite{leitnerRacialBiasAssociated2016}.} I ground my analysis in the theoretical framework of \textcite{akerlofEconomicsIdentity2000}, explicitly modeling how external prejudice creates differential utility from identity choices and establishing conditions under which individuals strategically modify their racial self-presentation. I also use CPS data to study the effect of education (parental and individual) and family income on Asian racial identity reporting.

Measuring identity choices outside of lab settings is challenging, requiring both objective ancestry indicators and subjective identity measures. I leverage birthplace and ancestry data to construct objective Asian ancestry measures, then analyze deviations between objective ancestry and subjective racial identity. I find that racial identity reporting responds to both individual characteristics and environmental factors reflecting local discrimination levels. 

I find substantial Asian racial identity attrition across generations, driven primarily by interracial family structures. While children of endogamous Asian parents report Asian identity at rates above 90\%, only about a third of second-generation children with one Asian and one White parent do the same. Using multinomial logit models, I show that anti-Asian bias shifts identity choices away from ``Asian only'' and toward ``White only,'' with effects that are larger in families where the father is Asian and that scale with the number of Asian grandparents among the third generation. Parental education promotes Asian identification on average, but this protective effect is undermined in high-bias states: education amplifies rather than buffers the identity-shifting effect of bias. The resulting negative selection on socioeconomic status among the Asian-identified population implies that conventional analyses overestimate Asian--White gaps precisely where prejudice is strongest.

This research contributes to several bodies of literature. First, it extends the economics of identity framework \autocite{akerlofEconomicsIdentity2000} and stratification economics \autocite{darityEconomicsIdentityOrigin2006,darityPositionPossessionsStratification2022} by examining how racial stereotypes---both positive and negative---influence identity choices among Asian Americans, who face a unique utility landscape where identity can simultaneously signal competence and foreignness \autocite{charnessSocialIdentityGroup2020,atkinHowWeChoose2021,goldsmithDarkLightSkin2007,hamiltonSheddingLightMarriage2009,dietteSkinShadeStratification2015}. Qualitative research documents how individuals strategically adapt racial presentations across contexts \autocite{waters1990ethnic,telles2008generations,zhou1997segmented,cornell2006ethnicity}; this paper quantifies those choices by comparing objective ancestry with subjective identification across varying bias environments.

Second, it extends research on discrimination \autocite{bertrandAreEmilyGreg2004,charlesPrejudiceWagesEmpirical2008,bursztynImmigrantNextDoor2022,bordaloStereotypes2016} by examining how prejudice influences racial self-identification itself. Within immigration research, the paper builds on studies of assimilation \autocite{abramitzkyCulturalAssimilationAge2016,abramitzkyNationImmigrantsAssimilation2014}, but unlike European immigrant groups, Asian Americans face persistent ``perpetual foreigner'' stereotypes that complicate integration regardless of generational status \autocite{foukaImmigrantsAmericansRace2022}, and the model minority myth means Asian racial identity may carry both benefits and costs depending on context \autocite{mengIntermarriageEconomicAssimilation2005}.

This paper most closely relates to recent economic research on racial identity fluidity \autocite{hadah2024hispanicidentity, antmanEthnicAttritionObserved2016,antmanIncentivesIdentifyRacial2015,antmanAmericanIndianCasinos2021}. While previous work focused primarily on Hispanic ethnic attrition, Asian American identity choices operate through different mechanisms due to distinct stereotypes, discrimination patterns, and socioeconomic profiles.

\section{Theoretical Framework}\label{sec:model}

I develop a theoretical framework for understanding racial identity choice that extends \textcite{akerlofEconomicsIdentity2000} to incorporate stereotype-specific costs and benefits. Unlike generic minority identification models, this framework recognizes that Asian Americans face unique utility trade-offs where racial identity can signal both positive attributes (academic achievement, work ethic) and negative characteristics (foreignness, social exclusion).

Formally, individual $i$ belongs to racial group $r_i \in \{A, W\}$, where $A$ represents Asian and $W$ represents White. Agent $i$'s utility depends on their actions and how those actions interact with their chosen racial identity $I_i$:

\begin{equation}
U_i = U_i(\pmb{a_i}, \pmb{a_{-i}}, I_i)\label{eq:util}
\end{equation}

Individual identity $I_i$ reflects personal actions, others' behaviors toward them, and societal expectations associated with their racial group:

\begin{equation}
I_i = I_i(\pmb{a_i}, \pmb{a_{-i}}; \pmb{S}_{r_{i}})\label{eq:identity}
\end{equation}

where $\pmb{a_i}$ represents individual $i$'s actions, $\pmb{a_{-i}}$ captures others' actions affecting $i$'s identity (including anti-Asian bias), and $\pmb{S}_{r_{i}}$ denotes societal stereotypes and expectations associated with racial group membership.\footnote{This extends \textcite{akerlofEconomicsIdentity2000}'s prescription concept to encompass both negative stereotypes and positive model minority expectations.}

The key insight is that $\pmb{S}_{A}$ includes both positive stereotypes (academic excellence) and negative ones (perpetual foreigner status). Thus creating context-dependent utility from Asian racial identification.

Individual $i$ selects actions $\pmb{a_i}$ to maximize utility given their racial group $r_i$, associated stereotypes $\pmb{S}_{r_{i}}$, and others' actions $\pmb{a_{-i}}$. The first-order condition becomes:

\begin{equation}
\frac{\partial U_i}{\partial \pmb{a_i}} + \frac{\partial U_i}{\partial I_i} \cdot \frac{d I_i}{d \pmb{a_i}} = 0\label{eq:foc}
\end{equation}

The solution $a_{i}^{\star}$ yields utility $U_{i}^{\star}$. Now suppose individuals can strategically choose their racial identity at cost $c$. They will switch identities when $\tilde{U_{i}}^{\star} \geq U_{i}^{\star} + c$, where $\tilde{U_{i}}^{\star}$ represents utility under the alternative racial identity.

Identity switching occurs when benefits $\tilde{U_{i}}^{\star} - U_{i}^{\star}$ exceed costs $c$. These net benefits are non-zero only when $\frac{d I_i}{d \pmb{a_i}} \neq 0$ and $\frac{\partial U_i}{\partial I_i} \neq 0$. This framework suggests empirical analysis should focus on: (1) individual characteristics affecting optimal actions under different racial identities, (2) contextual factors (anti-Asian bias) creating differential treatment by racial group, (3) populations with low switching costs $c$, and (4) groups whose utility significantly depends on racial identity.

The model predicts that anti-Asian bias increases the utility differential between White and Asian racial identity, making identity switching more attractive. Mixed-race individuals face lower switching costs due to phenotypic ambiguity, while later-generation Asian Americans may find identity switching more feasible due to cultural assimilation.

\section{Data Sources and Measurement Strategy}\label{sec:data}

In this section, I describe the datasets I use in the analysis. To examine relationships between social attitudes and Asian racial identity reporting, I need both subjective and objective Asian identity measures for identifying appropriate Asian American subgroups. I use the Integrated Public Use Microdata Series (IPUMS) Current Population Survey (CPS) \autocite{floodsarahIntegratedPublicUse2021a} with ancestry information through the places of birth of individuals, their parents, and grandparents to construct objective identity measures. I develop composite anti-Asian bias measures using \textcite{lubotskyInterpretationRegressionsMultiple2006}'s multiple proxy regression method to reduce attenuation bias.

\subsection{Measuring Asian Racial Identity}\label{subsec:cps}

I measure Asian racial identity using CPS data from 2004--2021. Following \textcite{antmanEthnicAttritionObserved2016,antmanEthnicAttritionAssimilation2020}, I utilize birthplace information for individuals, parents, and grandparents to create objective Asian ancestry indicators for minors under age 17 living with parents.\footnote{This approach parallels previous research but focuses on racial rather than ethnic categorization.} The methodology allows for the identification of first-, second-, and third-generation Asian Americans (see Figure \ref{fig:diag} for visual representation). I also use a CPS sample of adults aged 18+ to examine whether their racial identity reporting patterns differ from minors; for the adult sample, I can only identify first- and second-generation.\footnote{Since adults still living with parents might not be representative of the overall adult population and is a rare occurrence, I do not link adults to their parents.}

The objective ancestry measure---distinct from subjective racial identification where respondents select ``Asian'' as their race---depends on birthplaces across three generations.\footnote{For this analysis, Asian countries comprise East Asian and Southeast Asian nations, including China, Hong Kong, Taiwan, Japan, Korea, Mongolia, Cambodia, Indonesia, Laos, Malaysia, Philippines, Singapore, Thailand, and Vietnam, but exclude South Asian and Middle Eastern countries, consistent with standard demographic classifications.} The three identifiable generations include: 1) first-generation immigrants born in Asian countries with both parents also born in Asian countries, 2) second-generation individuals who are US-born citizens with at least one parent born in an Asian country, and 3) third-generation Asian Americans who are US-born citizens with two US-born parents and at least one grandparent born in an Asian country. 

While the ancestry measure provides an objective assessment of Asian heritage, it may not capture all nuances of racial identity. For instance, White individuals with Asian ancestry born to non-Asian parents in Asia, such as on American military bases, may be classified as Asian in the data. To avoid potential misclassification, I remove individuals who report that they were born abroad of American parents. The final sample includes Asian Americans, first-, second-, and third-generation immigrants aged 17 and younger living with parents between 2004 and 2021. I present the summary statistics in Table (\ref{tab:sumstat1}). The adult sample includes first- and second-generation Asian American immigrants aged 18 and older between 2004 and 2021. I show the summary statistics for the adult sample in Table (\ref{tab:sumstat-adults}).

CPS relies on household respondents (parents or caregivers) to report children's racial identity. Asian identity reporting rates are consistent regardless of whether mother (72\%), father (72\%), or child/other caregiver (87\%) serves as respondent (Table \ref{tab:hispbyproxy}), and attrition patterns among adults align with those observed in children (Tables \ref{tab:asianbygen} and \ref{tab:asianbygen-adults}). I address proxy reporting concerns more fully in Section~\ref{sec:robcheck}.

The overall sample comprises 49\% females, with 65\% self-reporting Asian racial identity---answering affirmatively to ``what is your race.'' Average age is 8.4 years. Approximately 52\% of mothers and 52\% of fathers hold college degrees. Additional summary statistics for the overall sample and each generation appear in Table (\ref{tab:sumstat1}).

Using parental and grandparental birthplaces, I objectively identify ethnic ancestry and categorize different family types. For second-generation children, parental birthplaces create three objective categories: Objectively Asian-father-Asian-mother (AA), Objectively Asian-father-White-mother (AW), and Objectively White-father-Asian-mother (WA)

Similarly, grandparental birthplaces create 15 objective categories for third-generation children: (1) objectively Asian paternal grandfather-Asian paternal grandmother-Asian maternal grandfather-Asian maternal grandmother (AAAA); (2) objectively White paternal grandfather-Asian paternal grandmother-Asian maternal grandfather-Asian maternal grandmother (WAAA); (3) objectively Asian paternal grandfather-White paternal grandmother-Asian maternal grandfather-Asian maternal grandmother (AWAA), etc.

My analysis employs a subsample of the US population; Tables \ref{tab:asianbygen} and \ref{tab:sumstat-adults} demonstrate sufficient observations across generations for both adults and children. Consistent with literature on ethnic and racial identity fluidity among Asian and Hispanic Americans, I document significant attrition among third-generation Asian Americans.\footnote{\textcite{duncanIdentifyingLaterGenerationDescendants2018,duncanSocioeconomicIntegrationImmigrant2018, antmanEthnicAttritionObserved2016,antmanEthnicAttritionAssimilation2020} document substantial identity attrition among various groups.}

\subsection{Measuring Anti-Asian Sentiment}\label{sub:lw-bias}

I construct anti-Asian sentiment measures using implicit association tests, American National Election Studies, and hate crimes targeting Asian Americans from 2004--2021. 

The implicit association test measures conceptual associations---for example, linking Asian Americans with negative stereotypes---and evaluative responses. Respondents rapidly categorize words into screen-displayed categories. Figure (\ref{fig:iatexamples}) shows examples from Harvard's Project Implicit skin tone test.

I employ Asian-focused IAT data to construct anti-Asian prejudice proxies \autocite{greenwaldMeasuringIndividualDifferences1998}. The IAT measures bias direction and magnitude while capturing unconscious biases individuals may be unwilling to report, and previous research demonstrates the difficulty of manipulating IAT scores \autocite{egloffPredictiveValidityImplicit2002}. A meta-analysis by \textcite{greenwaldMeasuringIndividualDifferences1998} finds significantly higher predictive validity for IAT compared to self-report measures.\footnote{IAT participation is voluntary, potentially creating selection bias. However, IAT-reflected bias serves as a proxy for prejudiced attitudes \autocite{chettyRaceEconomicOpportunity2020}.} 

However, some research questions IAT predictive validity claims. Implicit Association Tests may not reliably measure or predict implicit prejudice or biased behaviors. Research shows implicit biases experience minor, temporary intervention-induced changes. Additionally, implicit bias fails to predict dictator game contributions or social pressure susceptibility, highlighting distinctions between implicit bias and biased actions \autocite{arkesAttributionsImplicitPrejudice2004,forscherMetaanalysisProceduresChange2019,leeDoesImplicitBias2018}. Therefore, I supplement IAT with explicit bias measures from American National Election Studies (ANES) and hate crimes against Asian individuals to construct a composite bias measure.

I develop a second proxy from the American National Election Studies (ANES), averaging responses across several 2004--2020 questions measuring racial animus \autocite{anes2021}.\footnote{Questions parallel those used by \textcite{charlesPrejudiceWagesEmpirical2008}: (1) ``Conditions Make it Difficult for Blacks to Succeed'', (2) ``Blacks Should Not Have Special Favors to Succeed'', (3) ``Blacks Must Try Harder to Succeed'', (4) ``Blacks Gotten Less than They Deserve Over the Past Few Years'', and (5) ``Feeling Thermometer Toward Asians.''} While the ANES racial animus questions primarily focus on attitudes toward Black Americans, research demonstrates that racial prejudices are highly correlated across different minority groups, with individuals who express bias toward one racial minority typically holding similar attitudes toward others \autocite{almasalkhi2023links, mora2020antiblackness}. These measures therefore capture broader patterns of racial animus that extend beyond anti-Black sentiment specifically. When combined with the Asian-focused IAT measures and hate crimes against Asian Americans in my composite index, this multi-proxy approach weights the bias measure more heavily toward Asian-specific prejudice while still capturing the general racial climate. Finally, I incorporate Uniform Crime Reports (UCR) data quantifying hate crimes against Asian Americans \autocite{ucrbook}.

To reduce attenuation bias and measurement error, I follow \textcite{lubotskyInterpretationRegressionsMultiple2006} in constructing composite bias measures using IAT, ANES racial animus measures, and hate crimes against Asian Americans.\footnote{Additional methodological details appear in the Data Online Appendix, Section~\ref{sub:lw-bias}.} Figure (\ref{fig:skiniat}) graphically represents bias measures over time in the most and least biased locations. Figure (\ref{fig:Asian-twostates}) shows Asian racial identity reporting in the two most and least biased locations. Lower scores indicate less bias; higher scores indicate greater racial animus. One standard deviation bias increase is equivalent to moving from Washington, DC, or Vermont to North Dakota in 2020. State-level average bias over time appears in Figure (\ref{fig:skiniat-maps}), with overall 2004--2021 averages in Figure (\ref{fig:iat-map-all}).

\section{From the Data: Asian Racial Identity and Attrition}\label{sec:attrition}

Table (\ref{tab:asianbygen}) displays racial attrition patterns across generations. Among first-generation Asian Americans, 96\% self-report Asian racial identity. This drops to 73\% among second-generation and 31\% among third-generation Asian Americans. Attrition is driven primarily by children from interracial families. Among second-generation children, those with two Asian parents report 97\% Asian racial identity, while those with one Asian and one White parent report only 33\%. Similarly, among third-generation children, 94\% of those with four Asian grandparents report Asian racial identity compared to 25\% overall. Adult samples show similar patterns (Table \ref{tab:asianbygen-adults}). Among second-generation adults, those with two Asian parents report 95\% Asian racial identity, while those with one Asian and one White parent report 37\%.

Figures \ref{fig:histogram-all}--\ref{fig:histogram-thirdgen} display racial identity choices among objectively Asian children by generation and family structure.\footnote{Beginning in 2003, the CPS allowed respondents to report multiple racial identities. I categorize responses into six classifications: (1) Asian only, (2) White only, (3) Asian and White/Pacific Islander, (4) other non-Asian multiracial combinations, (5) Asian combined with other races, and (6) Asian/Pacific Islander.} Among all objectively Asian children (Figure \ref{fig:histogram-all}), 64\% report Asian only identity, 16\% White only, and 16\% Asian and White/Pacific Islander. First-generation children overwhelmingly report Asian only identity (94\%, Figure \ref{fig:histogram-firstgen}). Among second-generation interracial families, rates differ by parental composition: 36\% report Asian only in Asian father--White mother families versus 29\% in White father--Asian mother families (Figure \ref{fig:histogram-secondgen}).

Third-generation patterns reveal increasing identity fluidity (Figure \ref{fig:histogram-thirdgen}). White only identity becomes most common overall (35\%), followed by Asian only (29\%), Asian and White/Pacific Islander (26\%), and other multiracial identities (7\%). The number of Asian grandparents strongly predicts choices: those with one Asian grandparent report 54\% White only, while those with four Asian grandparents report 92\% Asian only.

These patterns demonstrate that racial identity becomes increasingly fluid across generations, with endogamous Asian families maintaining high rates of Asian racial identity while interracial family structures significantly increase not-Asian-only identity choices.

\section{Empirical Approach}\label{sec:empstrat}

To understand associations between Asian racial self-identification and anti-Asian bias, I estimate a multinomial logit model in which the dependent variable $Y_{ist}^g$ captures categorical racial identity choices: (1) Asian only, (2) White only, (3) Asian and White/Pacific Islander:

\begin{align}
\log\left(\frac{P(Y_{ist}^g = j)}{P(Y_{ist}^g = \text{Asian only})}\right) &= \beta_{1j}^g AntiAsianBias_{st} + \beta_{2j}^g DadCollegeGrad_{ist}\nonumber \\ 
& + \beta_{3j}^g MomCollegeGrad_{ist} + \beta_{4j}^g Female_{ist} + X_{ist}^g\pi_j  \nonumber \\ 
& + \gamma_{rtj} + \varepsilon_{istj}; \quad j \in \{\text{White only}, \text{Asian and White}\} \label{eq:multinomial_logit}
\end{align}

where $Y_{ist}$ represents the categorical racial identity of person $i$ in state $s$ at interview time $t$, $j$ indexes the identity categories, with ``Asian only'' as the reference category, $P(Y_{ist}^g = j)$ denotes the probability of person $i$ in state $s$ at time $t$ choosing identity $j$, $\mathbf{X}_{ist}^g$ represents the vector of controls that include quartic age, Asian population fraction, parent types, grandparent types, and generation dummies, $\boldsymbol{\beta}_j^g$ denotes the coefficient vector for outcome $j$ and generation $g$, and $J$ indexes the three identity categories. Finally, $\gamma_{rtj}$ represents region-time fixed effects that would control for region $\times$ year specific shocks affecting identity choice $j$, and $\varepsilon_{istj}$ is the error term.

The model specification allows the effects of anti-Asian bias and other covariates to vary across identity choices. For instance, $\beta_{1,\text{White only}}^g$ captures how anti-Asian bias affects the log-odds of choosing ``White only'' versus ``Asian only'' identity, while $\beta_{1,\text{Asian and White}}^g$ measures the bias effect on choosing ``Asian and White'' versus ``Asian only.'' This framework enables analysis of how bias, sex, and parental education differentially influence various identity strategies available to individuals with Asian ancestry.

The coefficients of interest are $\beta_{1j}^g$, $\beta_{2j}^g$, $\beta_{3j}^g$, and $\beta_{4j}^g$, which capture how anti-Asian bias, parental education, and sex influence the likelihood of selecting each identity category relative to ``Asian only.'' I estimate separate models for each generation $g \in \{1,2,3\}$ to assess whether these relationships differ across generational status.

The multinomial specification is appropriate because choosing ``White only'' identity represents a fundamentally different strategy than selecting ``Asian and White'' multiracial identity, a distinction particularly important for mixed-ancestry individuals who may strategically choose between complete ethnic distancing and maintaining partial ethnic connection.

Since log odds coefficients from multinomial logit models are difficult to interpret directly, I compute the predicted probabilities of each identity choice at different anti-Asian bias levels, holding other covariates constant and summarizing results with the median across 1{,}000 bootstrap resamples (with percentile-based 95\% confidence intervals). This approach provides more intuitive insights into how bias influences the likelihood of selecting each identity category.

\section{Results and Robustness Checks}\label{sec:results}

I present average marginal effects from the multinomial logit analysis in Figures \ref{fig:marginal-effects-second-pooled}--\ref{fig:marginal-effects-third-grandparental}, focusing on mixed-race Asian Americans where effects are most pronounced.

\textbf{Second-Generation Subsample.} Among children of mixed-race parents (Figure \ref{fig:marginal-effects-second-pooled}), a one standard deviation increase in anti-Asian bias increases ``White only'' identity by 12 percentage points (pp) and decreases ``Asian only'' by 10~pp. Maternal college education increases ``Asian only'' ($+$5~pp) and decreases ``White only'' ($-$6~pp). Paternal college education increases ``Asian and White'' ($+$5~pp) and decreases ``White only'' ($-$4~pp). Results among second-generation adults are directionally similar (Figure \ref{fig:marginal-effects-second-adults-pooled}).

I also estimate the model separately by parental composition---Asian father--White mother (AW) versus White father--Asian mother (WA) families (Figure~\ref{fig:marginal-effects-second-parental} and Table \ref{tab:between-model-second}). Point estimates suggest that bias effects on ``White only'' identity may be larger in AW families ($+$20~pp) than WA families ($+$10~pp), and that maternal college education effects may differ in direction and magnitude across the two family types. I conduct a formal bootstrapped between-model tests of cross-group differences in average marginal effects (Figure~\ref{fig:between-model-second} and Table \ref{tab:between-model-second}).

I find that the AW--WA difference in bias effects on ``White only'' identity is statistically significant ($p < 0.05$). Consequently, the effect of bias is stronger among AW families (when the father is Asian), pushing children toward ``White-only'' racial identity. Maternal and paternal college education effects also differ significantly across family types for ``Asian only'' and ``White only'' outcomes (all $p < 0.05$). However, I find that the remaining AW--WA differences are not statistically distinguishable. 

\textbf{Third-Generation Subsample.} Among pooled third-generation children, bias effects on identity are not statistically significant. Paternal college education increases ``Asian and White'' ($+$8~pp). I find that bias and parental education have larger effects on racial identity among those who have more Asian grandparents, so I split the sample by number of Asian grandparents (Figure~\ref{fig:marginal-effects-third-grandparental}). The effect of bias on ``White only'' identity rises from near zero (one Asian grandparent) to $+$19~pp (three Asian grandparents), and maternal college effects on ``Asian only'' rise from near zero to $+$32~pp. 

Bootstrapped between-model tests (Figure~\ref{fig:between-model-third} and Table \ref{tab:between-model-third}) confirm that some, but not all, of these apparent differences are statistically significant. Comparing one versus three Asian grandparents---the most extreme contrast---the bias effect on ``White only'' differs significantly ($p < 0.01$), as do maternal and paternal college effects on ``Asian only'' and ``White only'' (all $p < 0.05$ or better). Comparing one versus two grandparents, maternal college effects differ significantly on ``White only'' ($p < 0.05$) and ``Asian and White'' ($p < 0.01$), but bias and paternal college effects are not distinguishable. Comparing two versus three grandparents, bias effects on ``White only'' and parental education effects on ``Asian only'' and ``White only'' all differ significantly ($p < 0.01$), but gender effects do not.

I find that the effects of bias and parental education on identity are heterogeneous across the different ancestry types. Even where cross-group differences are not statistically significant, I find that the point estimates grow larger as the number of Asian grandparents increases. Consistent with the theoretical framework, individuals with lower switching costs---those with fewer Asian grandparents and greater racial ambiguity---show smaller marginal effects of bias because they can switch identities even at low levels of bias, while those with more Asian ancestry require stronger incentives to overcome higher switching costs.\footnote{\textcite{hadah2024hispanicidentity} finds similar effects on third-generation Hispanic immigrants.}
\subsection*{The interaction between state-level bias and socioeconomic status} To assess whether bias and socioeconomic status jointly affect Asian identification, I estimate linear probability models interacting demeaned state-level bias with parental education (own education for adults), family income (own income for adults), and gender. Results appear in Figures~\ref{fig:interaction-coefs-aw-wa-awwa}--\ref{fig:interaction-coefs-thirdgen-grandparent}.
 
The central finding is that education amplifies the negative effect of bias on Asian identification. Among second-generation AW children (Figure \ref{fig:interaction-coefs-aw-wa-awwa}), the interaction between bias and maternal-college is negative and significant (approximately $-$0.20); in high-bias states, children of college-educated mothers are substantially less likely to report ``Asian only'' identity. AW adults show the same pattern with own college education (approximately $-$0.07; Figure~\ref{fig:interaction-coefs-aw-wa}). Among WA families, no interaction terms are significant in either the children or adult samples. In the pooled second-generation sample, the bias-by-maternal-college interaction is negative and marginally significant (approximately $-$0.04). Among third-generation respondents (Figure~\ref{fig:interaction-coefs-thirdgen-grandparent}), point estimates follow the same direction---negative bias-by-maternal-college interactions that grow larger with more Asian grandparents---but are imprecisely estimated. The pooled third-generation bias-by-maternal-college interaction is significant (approximately $-$0.05).
 
These results imply that higher-educated individuals of Asian ancestry are disproportionately likely to exit Asian identification in high-bias states. The remaining Asian-identified population is therefore negatively selected on education and socioeconomic status, implying that conventional analyses could overestimate Asian--White gaps in states where anti-Asian prejudice is strongest.

\subsection{Robustness Checks}\label{sec:robcheck}

I examine three potential threats to identification: endogenous interracial marriage, selective migration, and proxy reporting. If anti-Asian bias influences interracial marriage rates, it could affect the composition of the mixed-race sample. I regress interracial couple status on state-level bias with partner-specific controls (Table~\ref{regtab-logit-02}). A one standard deviation increase in bias raises interracial parent probabilities by 4~pp, driven by a 3~pp increase in Asian wife--White husband couples and a 1~pp decrease in Asian husband--White wife couples.

Using 2004--2021 Census data \autocite{floodsarahIntegratedPublicUse2021}, I test whether anti-Asian bias predicts interstate migration among second-generation Asian Americans with two Asian-born parents (Table~\ref{regtab-mig-01}). I find no significant correlation between bias and migration decisions. Among movers, Asian-identifying individuals live in states with 0.06 standard deviations \emph{greater} bias than their birth states, suggesting that if anything my estimates understate the effect of bias on identity switching.

In the CPS, a household respondent reports racial identity for all members, raising the concern that proxy responses may not reflect children's own identities. Two pieces of evidence mitigate this concern. First, \textcite{duncanIntermarriageIntergenerationalTransmission2011} show that reported Hispanic identity does not vary with respondent identity, and Table~\ref{tab:hispbyproxy} shows a similar pattern for Asian identity.\footnote{According to CPS guidelines, household respondents must be at least 15 years old with sufficient household knowledge. When the proxy is ``self,'' respondents range from 15 to 17 years old.} Second, I find similar patterns of ethnic attrition and bias effects among adult Asian Americans who self-report their own identity (Table~\ref{tab:asianbygen-adults}).

Larger Asian American populations could themselves affect state-level bias. Figure~\ref{scatter-plot-1} shows no correlation between state Asian American population shares and average anti-Asian bias.

\section{Conclusion}\label{sec:conc}

As American society becomes increasingly multiracial, racial identity choices will significantly influence political representation, resource allocation, and social cohesion. Understanding the determinants of identity is particularly important for researchers studying discrimination's role in racial economic gaps. This paper demonstrates how individual characteristics and anti-Asian sentiment influence racial identity reporting among Asian Americans.
 
Using multinomial logit analysis, I show that anti-Asian bias fundamentally reshapes racial identity choices, with effects that vary significantly across family structures and generational distance from immigration. Among pooled second-generation children of mixed-race parents, a one standard deviation increase in bias raises ``White only'' reporting by 12 percentage points and reduces ``Asian only'' by 10~pp. The effects differ markedly by parental composition: in AW families, the same increase in bias reduces ``Asian only'' by 13~pp and increases ``White only'' by 20~pp, while in WA families, ``Asian only'' decreases by 12~pp and ``White only'' increases by 10~pp. Between-model tests confirm that the bias effect on ``White only'' identity is significantly larger in AW families (10.6~pp, $p < 0.05$). Among third-generation respondents, bias effects scale with the number of Asian grandparents: those with three Asian grandparents show the largest shifts ($+$19~pp ``White only''), while those with one grandparent show small, imprecise effects. 
 
Parental education shapes identity choices in ways that also differ significantly across family structures. In AW families, maternal college education increases ``Asian only'' by 10~pp and decreases ``White only'' by 12~pp, while in WA families maternal education produces no significant change---a difference confirmed by formal between-model tests on both outcomes ($p < 0.05$). Among the third generation, parental education effects grow dramatically with ancestry concentration: maternal college raises ``Asian only'' by 32~pp for those with three Asian grandparents versus 7~pp for those with two, differences that are statistically significant ($p < 0.01$). However, the interaction between bias and education reveals that these protective effects are undermined in high-bias environments: the bias-by-maternal-college interaction is $-$0.20 in AW families, indicating that education amplifies rather than buffers the negative effect of bias on Asian identification.
 
These results have important consequences for the measurement of racial disparities. Because the combination of higher education and higher bias is associated with less---not more---Asian identification, the Asian-identified population in high-bias states is negatively selected on education and socioeconomic status. Conventional analyses relying on self-reported race therefore overestimate Asian--White gaps in the most prejudiced areas. More broadly, identity decisions likely profoundly affect people's choices, investments, and well-being, and the strategic nature of these decisions means that studies relying on self-reported racial identity may overstate Asian American socioeconomic success and understate assimilation speed.
 
This study encourages further research into relationships between bias and self-reported identities for other groups, including other ethnic and racial minorities and sexual minorities. Scholars could examine how recent anti-Asian violence following the COVID-19 pandemic has influenced identity patterns, providing natural experiments in bias effects. Researchers might also explore identity choices in specific institutional contexts like college admissions or workplace advancement, where model minority stereotypes create complex incentive structures, or extend the analysis to subgroups facing distinct stereotypes and discrimination patterns. More granular geographic analysis at the county or zip-code level could provide more precise measures of local bias, though data limitations prevent this in the current analysis. Understanding strategic racial identification among Asian Americans is essential for designing effective anti-discrimination policies and accurately measuring racial equity progress.
